\section{Selección del modelo final y despliegue}

\subsection{Criterio de selección}

\textbf{ConvNeXt-Small} se selecciona como modelo de producción principal por su mayor accuracy (96.47\%) y F1 macro (96.47\%), indicando rendimiento equilibrado en las 15 clases. La diferencia con el baseline de Swin-T (+5.4 pp) y la consistencia de las configuraciones top-4 (todas por encima del 95.8\%) refuerzan la confianza en su generalización.

\textbf{EfficientNetV2-S} constituye la alternativa recomendada cuando se requiere baja latencia: con un 96.00\% de accuracy y $2.3\times$ menor tiempo por época (43.5\,s vs.\ 99.9\,s), es la opción preferible en entornos con restricciones de cómputo, como instancias serverless o despliegues sin GPU.

\subsection{Arquitectura del sistema de inferencia}

El servicio de inferencia en producción está compuesto por dos componentes independientes comunicados a través de HTTP:

\begin{itemize}\setlength\itemsep{3pt}
    \item \textbf{Backend (FastAPI):} expone los siguientes endpoints:
    \begin{itemize}\setlength\itemsep{2pt}
        \item \texttt{GET /} — health check rápido; devuelve \texttt{status} y nombre del modelo cargado.
        \item \texttt{GET /health} — health check detallado; devuelve \texttt{status}, \texttt{model\_loaded} y \texttt{device} (CPU/GPU).
        \item \texttt{GET /classes} — devuelve la lista de las 15 categorías con su índice.
        \item \texttt{GET /model-info} — devuelve metadatos del modelo: arquitectura, ruta del checkpoint, tamaño de entrada, número de clases y dispositivo.
        \item \texttt{POST /predict} — recibe una imagen en base64 y devuelve la clase predicha, la confianza y el nombre de fichero.
        \item \texttt{POST /predict-topk} — recibe una imagen en base64 y el parámetro \texttt{k}; devuelve las \textit{k} clases más probables con sus confidencias, permitiendo al frontend mostrar la distribución completa de predicciones.
    \end{itemize}
    La documentación interactiva (Swagger UI) está disponible en \texttt{/docs} y la estática en \texttt{docs/api.md}.
    \item \textbf{Frontend (Streamlit):} interfaz web que permite cargar imágenes locales o por URL, muestra la predicción principal con su nivel de confianza y visualiza la distribución de probabilidad sobre las 15 clases mediante gráfico de barras interactivo.
\end{itemize}

\subsection{Rendimiento por clase y viabilidad operativa}

Al finalizar cada run, el script registra automáticamente en W\&B la matriz de confusión sobre el conjunto de validación, vinculada al checkpoint y configuración exactos de ese run. Esto permite un análisis de errores trazable y reproducible.

Del análisis de las matrices de confusión se desprenden patrones consistentes: \textit{Highway} y \textit{Suburb} alcanzan el 100\% en los tres modelos, mientras que \textit{Open country} es la clase más difícil (ConvNeXt-Small: 87\%, Swin-T: 75\%), confundiéndose principalmente con \textit{Coast} y \textit{Mountain}. La mayor diferencia entre modelos aparece en clases interiores: Swin-T clasifica \textit{Bedroom} con solo un 78\% (16/100 confundidas con \textit{Living room}), frente al 94\%+ de ConvNeXt-Small y EfficientNetV2-S, lo que explica la brecha de 4~pp en F1 macro.

El F1 macro de 96.47\% de ConvNeXt-Small supera ampliamente el umbral de éxito ($\geq 80\%$) y los errores residuales se concentran en pares de clases ambiguos incluso para un observador humano, lo que los hace aceptables para el caso de uso previsto.
